\section{Related Work}

\paragraph{Foundation Models for 3D Medical Segmentation.}
The Segment Anything Model (SAM)~\citep{kirillov2023sam} introduced promptable segmentation for natural images, demonstrating remarkable zero-shot generalization through point, box, and mask prompts. SAM2~\citep{ravi2024sam2} extended this paradigm to video with memory-augmented tracking, enabling temporal consistency across frames. The adaptation of these models to 3D medical imaging has become a major research direction.

SAM-Med3D~\citep{wang2023sam_med3d} trains a volumetric SAM architecture on large-scale CT data with point prompts, achieving competitive performance on multi-organ segmentation. MedSAM~\citep{ma2024medsam} fine-tunes the original SAM for medical bounding-box prompts across diverse imaging modalities. Medical SAM2~\citep{zhu2024medsam2} reframes 3D medical segmentation as video object tracking via SAM2, introducing a self-sorting memory bank for cross-slice consistency. MA-SAM~\citep{chen2024masam} proposes modality-agnostic adapters that enable cross-modal transfer without modality-specific training. SAM2-Adapter~\citep{chen2024sam2adapter} introduces task-specific adapters to overcome SAM2 limitations in medical, camouflage, and shadow segmentation domains.

A common thread across these works is evaluation via segmentation accuracy metrics (Dice, IoU) without assessing whether the model faithfully follows prompt targets. Our work reveals that high Dice scores can coexist with catastrophic prompt unfaithfulness---a gap that standard evaluation completely misses.

\paragraph{Parameter-Efficient Adaptation for Medical Imaging.}
LoRA~\citep{hu2022lora} enables efficient fine-tuning by injecting trainable low-rank matrices into frozen model layers, dramatically reducing the number of trainable parameters while maintaining adaptation quality. This approach has been widely adopted in medical imaging.

LoRA-MedSAM~\citep{wu2024loramedsam} applies LoRA to MedSAM, reducing fine-tuning time by 22.7\% while improving Dice scores. FATE-SAM~\citep{deng2025fatesam} achieves few-shot 3D segmentation by reassembling pre-trained SAM2 modules without any gradient-based training, leveraging support examples for prompt-free inference. tCURLoRA~\citep{cheng2025tcurlora} proposes tensor CUR decomposition as an alternative to standard LoRA, achieving better parameter efficiency for medical segmentation tasks.

Our work uses standard LoRA adaptation but reveals a crucial insight: the \emph{training paradigm} (paired vs.\ single prompt) matters far more than the adaptation architecture or parameter count. Even with identical LoRA configurations, single-prompt training produces unfaithful models while paired-prompt training produces faithful ones.

\paragraph{Prompt Robustness in Medical Segmentation.}
Recent work has begun examining how medical segmentation models respond to imperfect prompts. \citet{chattopadhyay2026robustness} systematically evaluate how foundational 3D medical segmentation models respond to controlled perturbations of dense visual prompts, revealing that models rely heavily on spatial shape cues and exhibit varying resilience to different perturbation types. \citet{huang2024perturbed} study perturbed prompts for robust SAM adaptation in polyp segmentation, proposing augmentation strategies that improve stability. \citet{li2024saliency} develop saliency-guided prompt distillation to maintain segmentation quality under noisy or generic prompts in clinical workflows.

These works study \emph{robustness}---stability of predictions under small input perturbations. We identify a fundamentally different and more severe failure: \emph{faithfulness}---whether the model responds to target changes at all. The distinction is critical: a robust-but-unfaithful model produces stable predictions that consistently ignore the user, while a faithful model correctly changes its output when the user indicates a different target. SAM-Med3D is robust (1.3pt Dice drop at 10-voxel shift) but unfaithful (Switch Accuracy 20.8\%)---demonstrating that these are orthogonal properties requiring fundamentally different solutions.

\paragraph{Contrastive and Paired Training in Medical Analysis.}
Contrastive learning has been extensively applied to medical image analysis for representation learning and pre-training~\citep{zhao2024contrastive_survey}. Recent advances address specific challenges: \citet{liu2024collapse} identify and overcome dimensional collapse in self-supervised contrastive learning for medical images, showing that high inter-image similarity in medical data causes representation degeneration. \citet{wang2025domain_contrastive} propose domain and task-focused example selection for data-efficient contrastive medical segmentation.

Our paired prompt training shares the principle of exposing the model to contrasting conditions within a single optimization step. However, it differs fundamentally in objective: rather than learning invariant or discriminative \emph{representations}, we enforce prompt-conditioned output \emph{differentiation}---the model must produce different segmentation masks for different prompts. This is a supervised, task-level constraint rather than a representation-level one, and it directly addresses the prompt faithfulness failure rather than general feature quality.
